\subsection*{边界向右移动，只有一条窄带发生变化}
定义 $F(x)$ 为从固定起点累积到 $x$ 的面积。比较 $F(x+h)$ 与 $F(x)$，大部分旧区域完全相同，相减后只剩边界附近的一条窄带。因此求面积函数的变化率，不需要重新计算整块面积。

新增面积除以宽度 $h$，就是这条带的平均高度。如果原函数在边界连续，带内所有高度在足够窄时都接近 $f(x)$，平均高度也被夹在附近，所以极限就是 $f(x)$。连续性在这里负责让“小区间里高度差不多”成为可靠判断。

\subsection*{把近似拆成可看见的两部分}
取 $f(t)=t$，从 $x$ 到 $x+h$ 的面积包含高度为 $x$ 的矩形和高为 $h$ 的小三角形。总增量 $xh+h^2/2$ 除以 $h$ 后为 $x+h/2$，趋向边界高度。三角形就是有限步长近似的误差来源。

注意 $F(x)$ 的高度与 $F'(x)$ 的斜率不是同一种量：前者已经累积，后者描述新贡献进来的速度。某个点两者数值碰巧相同，不意味着两函数相同；检查单位也能发现这个区别。

反向移动边界时，窄带被减去，宽度也带负号，同样得到一致的导数。由此可见有向积分不是人为增加的麻烦，而是让累积与变化在两个方向都遵守同一规则。下一节将继续追问：变化速度自己又如何改变？
